% Level 1, Workbook chapter 11 -- Project Studio   (input by student-workbook.tex)
\chapterweek{11}{Project Studio}{Build, test and improve your project in small steps.}

\section{What you will learn}
\begin{itemize}[leftmargin=1.2em,itemsep=1pt]
  \item How real programmers build big programs: one small step at a time.
  \item How to test a partner's program and give helpful feedback.
  \item How to finish Milestones M2 and M3 of your art show.
\end{itemize}

\section{The big idea: build, run, check}
Big programs are not written in one go. Programmers go round this loop all day long:
\begin{center}
%% alt: A cycle of three steps. Build one small thing. Run it by pressing F5. Check: did it do what I wanted? Then build again. If something breaks, you know it was the last small change.
\begin{tikzpicture}[font=\small,
  n/.style={draw,thick,circle,minimum size=2.2cm,align=center}]
  \node[n,fill=kcsky!25] (b) at (90:2.1) {\textbf{Build}\\one small\\thing};
  \node[n,fill=kcyellow!35] (r) at (-30:2.1) {\textbf{Run}\\press F5};
  \node[n,fill=kcgreen!20] (c) at (210:2.1) {\textbf{Check}\\Did it do\\what I wanted?};
  \draw[-{Stealth},very thick] (b) to[bend left=25] (r);
  \draw[-{Stealth},very thick] (r) to[bend left=25] (c);
  \draw[-{Stealth},very thick] (c) to[bend left=25] (b);
  \node[align=left,anchor=west] at (3.6,0) {If something breaks,\\you know it was the\\\textbf{last small change}.};
\end{tikzpicture}
\end{center}
It is like building a LEGO set bag by bag: finish bag 1, check it looks like the picture, then open bag 2.

\subsection{Your milestones}
\begin{center}
%% alt: Four project milestones in a row. M1 Plan, in week 10: the plan sheet. M2 Core, in week 11 part 1: my function draws. M3 Surprise, in week 11 part 2: a question, an if and something random. M4 Showcase, in week 12: run and explain.
\begin{tikzpicture}[font=\small,
  m/.style={draw,thick,rounded corners=4pt,minimum width=3cm,minimum height=1.4cm,align=center}]
  \node[m,fill=kcgray!15] (m1) at (0,0) {\textbf{M1 Plan}\\week 10\\\faCheck\ plan sheet};
  \node[m,fill=kcsky!25] (m2) at (3.8,0) {\textbf{M2 Core}\\week 11, part 1\\my function draws};
  \node[m,fill=kcorange!25] (m3) at (7.6,0) {\textbf{M3 Surprise}\\week 11, part 2\\question, if, random};
  \node[m,fill=kcgreen!20] (m4) at (11.4,0) {\textbf{M4 Showcase}\\week 12\\run and explain};
  \draw[-{Stealth},very thick] (m1) -- (m2);
  \draw[-{Stealth},very thick] (m2) -- (m3);
  \draw[-{Stealth},very thick] (m3) -- (m4);
\end{tikzpicture}
\end{center}

\subsection{Five small steps (the tiny art show from the demo)}
\begin{enumerate}[leftmargin=1.4em,itemsep=1pt]
  \item \textbf{Question}: \texttt{how\_many = int(input("How many bubbles? "))}. Run.
  \item \textbf{Function}: \texttt{def bubble(x, y, color):} and call it once. Run.
  \item \textbf{Loop}: put the call inside \texttt{for i in range(how\_many):}. Run.
  \item \textbf{Random}: random \texttt{x}, \texttt{y} and color. Run twice!
  \item \textbf{If}: if the number is too big, use a smaller one. Run with 500.
\end{enumerate}
\begin{tip}[Save a version]
Before a big change, click \menu{File} → \menu{Save as…} and add \texttt{\_v2} to the name (\texttt{fish\_party\_v2.py}). If the change goes wrong, version 1 is still safe.
\end{tip}

\begin{exercise}[Sprint 1: core drawing works (M2)]
Open your project file. Change the example \texttt{my\_shape} function into \textbf{your} function: give it a name that says what it draws, and change the drawing. Draw it once at \texttt{(0, 0)}. Run. Then call it a second time somewhere else. Run.
\begin{checklist}
  \item My function has a good name: \rule{5cm}{0.4pt}
  \item It draws my main thing.
  \item I call it at least twice, in different places.
  \item The program runs with no red errors.
\end{checklist}
\end{exercise}

\begin{exercise}[Partner bug hunt]
Swap seats with your partner. Run their program \textbf{twice}: once with a normal answer, once with a strange one (0, a giant number, a word, a capital letter). Write what happened --- don't fix it, only report it!
\begin{center}
\renewcommand{\arraystretch}{1.6}
\begin{tabular}{@{}p{4cm}p{9.4cm}@{}}
\toprule
I typed & What happened\\
\midrule
 & \\ \midrule
 & \\
\bottomrule
\end{tabular}
\end{center}
Then give \textbf{two stars and a wish}:\par
\faStar\ I liked \rule{11.5cm}{0.4pt}\par\smallskip
\faStar\ I liked \rule{11.5cm}{0.4pt}\par\smallskip
\faMagic\ I wish \rule{11.5cm}{0.4pt}
\end{exercise}

\begin{exercise}[Sprint 2: surprise and polish (M3)]
Add the rest of your six ingredients. Check off each one in the comment list at the top of your file \textbf{when you have seen it work}.
\begin{checklist}
  \item My question(s) work (numbers use \texttt{int(...)}).
  \item My \texttt{if} changes something you can see.
  \item My loop repeats something: my function, or stars, bubbles, rays\ldots
  \item \texttt{random} makes it different every run.
  \item A title with \texttt{t.write(...)} and my \textbf{nickname} as the artist (Part 7 of the starter shows how).
\end{checklist}
\end{exercise}

\begin{tip}[Feedback is a gift]
\textbf{Giving}: be kind and specific (\enquote{I like how the stars are different sizes}), and make wishes about the program, never the person. \textbf{Getting}: say thank you --- then \textbf{you} decide what to change. It is your art.
\end{tip}

\begin{vocab}
  \vocabword{Iterate}{To improve something by going round the build--run--check loop again and again.}
  \vocabword{Milestone}{A checkpoint on the way to finishing a project.}
  \vocabword{Playtest}{Letting someone else try your program to find problems.}
  \vocabword{Feedback}{What someone tells you about your work, to help you improve it.}
  \vocabword{Version}{A saved copy of your program at one moment, like \texttt{\_v2}.}
\end{vocab}

\begin{reflection}
What bug did you find today? How did you find it?
\answerlines{2}
What feedback did you get? What will you change because of it?
\answerlines{2}
\end{reflection}
